\newcommand{\coursemodule}{Module 07}
% ════════════════════════════════════════════════════════════════════════════
% Pixels, Frames & Transformers: shared Beamer preamble (WUR style)
% Adapted from the preamble of Bags, Vectors & Transformers.
%
% Use in a module deck (from <part>/<module>/slides/):
%   \newcommand{\coursemodule}{Module 01}
%   % ════════════════════════════════════════════════════════════════════════════
% Pixels, Frames & Transformers: shared Beamer preamble (WUR style)
% Adapted from the preamble of Bags, Vectors & Transformers.
%
% Use in a module deck (from <part>/<module>/slides/):
%   \newcommand{\coursemodule}{Module 01}
%   % ════════════════════════════════════════════════════════════════════════════
% Pixels, Frames & Transformers: shared Beamer preamble (WUR style)
% Adapted from the preamble of Bags, Vectors & Transformers.
%
% Use in a module deck (from <part>/<module>/slides/):
%   \newcommand{\coursemodule}{Module 01}
%   \input{../../../shared/preamble.tex}
%   \begin{document}
%   \moduletitle{Images as Data}
%   ...
% Compile with: latexmk -pdf <deck>.tex
% ════════════════════════════════════════════════════════════════════════════

\documentclass[aspectratio=169]{beamer}

% ── Packages ────────────────────────────────────────────────────────────────
\usepackage[T1]{fontenc}
\usepackage[utf8]{inputenc}
\usepackage{lmodern}
\usepackage{microtype}
\usepackage{xcolor}
\usepackage{tikz}
\usetikzlibrary{arrows.meta, positioning, calc, shapes.geometric, decorations.pathreplacing, matrix, fit}
\usepackage{booktabs}
\usepackage{fontawesome5}
\usepackage{array}
\usepackage[most]{tcolorbox}

\providecommand{\coursemodule}{}

% ── WUR color palette ───────────────────────────────────────────────────────
\definecolor{wurgreen}{HTML}{34B233}
\definecolor{darkgreen}{HTML}{1F6B1E}
\definecolor{lightgreen}{HTML}{E8F7E8}
\definecolor{darknavy}{HTML}{1A1A2E}
\definecolor{midgray}{HTML}{6B7280}
\definecolor{lightgray}{HTML}{F3F4F6}
\definecolor{bodytext}{HTML}{1F2937}
\definecolor{codegreen}{HTML}{9AE6B4}
\definecolor{warnred}{HTML}{FEE2E2}
\definecolor{warnborder}{HTML}{EF4444}

% ── Beamer theme ────────────────────────────────────────────────────────────
\usetheme{default}
\usecolortheme{default}
\usefonttheme{professionalfonts}
\setbeamertemplate{navigation symbols}{}
\setbeamertemplate{itemize items}{\textcolor{wurgreen}{\textbullet}}
\setbeamertemplate{itemize subitem}{\textcolor{midgray}{--}}
\setbeamertemplate{enumerate items}{\textcolor{wurgreen}{\textbf{\insertenumlabel.}}}
\setbeamercolor{frametitle}{fg=white, bg=darknavy}
\setbeamertemplate{frametitle}{%
  \begin{beamercolorbox}[wd=\paperwidth, ht=1.1cm, dp=0.2cm, leftskip=0.6cm]{frametitle}
    \usebeamerfont{frametitle}\insertframetitle
  \end{beamercolorbox}%
}
\setbeamerfont{frametitle}{size=\large, series=\bfseries}
\setbeamercolor{background canvas}{bg=white}
\setbeamercolor{normal text}{fg=bodytext}
\setbeamerfont{normal text}{size=\small}
\setbeamertemplate{footline}{%
  \begin{beamercolorbox}[wd=\paperwidth, ht=0.4cm, dp=0.15cm]{footline}%
    \hspace{0.5cm}%
    \textcolor{midgray}{\tiny Pixels, Frames \& Transformers · \coursemodule\ · Wageningen University \& Research}%
    \hfill%
    \textcolor{midgray}{\tiny \insertframenumber}%
    \hspace{0.4cm}%
  \end{beamercolorbox}%
}

% ── Boxes ───────────────────────────────────────────────────────────────────
\tcbset{wurbase/.style={arc=3pt, boxrule=0.6pt, left=7pt, right=7pt, top=5pt, bottom=5pt,
                        fontupper=\small, before skip=4pt, after skip=4pt}}
\newtcolorbox{greenbox}[1][]{wurbase, colback=lightgreen, colframe=wurgreen, #1}
\newtcolorbox{graybox}[1][]{wurbase, colback=lightgray, colframe=lightgray, #1}
\newtcolorbox{navybox}[1][]{wurbase, colback=darknavy, colframe=darknavy, colupper=white, #1}
\newtcolorbox{warnbox}[1][]{wurbase, colback=warnred, colframe=warnborder, #1}
\newtcolorbox{codebox}[1][]{wurbase, colback=darknavy, colframe=darknavy, colupper=white,
                            fontupper=\ttfamily\scriptsize, left=6pt, top=4pt, bottom=4pt, #1}

% ── Title slide ─────────────────────────────────────────────────────────────
\newcommand{\moduletitle}[1]{%
{
\setbeamertemplate{footline}{}
\begin{frame}[plain]
  \begin{tikzpicture}[remember picture, overlay]
    \fill[darknavy] (current page.south west) rectangle (current page.north east);
    \fill[wurgreen] (current page.south west) rectangle ([xshift=0.45cm]current page.north west);
    \node[anchor=west, text=white, font=\huge\bfseries, text width=13cm]
      at ([xshift=1.1cm, yshift=1.2cm]current page.west) {Pixels, Frames \& Transformers};
    \node[anchor=west, text=wurgreen, font=\large, text width=13cm]
      at ([xshift=1.1cm, yshift=0.05cm]current page.west) {#1 · \coursemodule};
    \node[anchor=west, text=midgray, font=\small, text width=12cm]
      at ([xshift=1.1cm, yshift=-1.5cm]current page.west)
      {Denise J. Roth, PhD Candidate\\
       Strategic Communication Group · Wageningen University \& Research};
  \end{tikzpicture}
\end{frame}
}}

% ── Small helpers ───────────────────────────────────────────────────────────
\newcommand{\keyterm}[1]{\textbf{\textcolor{darkgreen}{#1}}}
\newcommand{\muted}[1]{\textcolor{midgray}{#1}}
\newcommand{\credit}[1]{{\tiny\textcolor{midgray}{#1}}}
\newcommand{\notebookslide}[1]{%
\begin{frame}{In the Notebook}
\vspace{0.1cm}
#1
\end{frame}}
\newcommand{\nbitem}[2]{\item[\textcolor{wurgreen}{#1}] #2}

%   \begin{document}
%   \moduletitle{Images as Data}
%   ...
% Compile with: latexmk -pdf <deck>.tex
% ════════════════════════════════════════════════════════════════════════════

\documentclass[aspectratio=169]{beamer}

% ── Packages ────────────────────────────────────────────────────────────────
\usepackage[T1]{fontenc}
\usepackage[utf8]{inputenc}
\usepackage{lmodern}
\usepackage{microtype}
\usepackage{xcolor}
\usepackage{tikz}
\usetikzlibrary{arrows.meta, positioning, calc, shapes.geometric, decorations.pathreplacing, matrix, fit}
\usepackage{booktabs}
\usepackage{fontawesome5}
\usepackage{array}
\usepackage[most]{tcolorbox}

\providecommand{\coursemodule}{}

% ── WUR color palette ───────────────────────────────────────────────────────
\definecolor{wurgreen}{HTML}{34B233}
\definecolor{darkgreen}{HTML}{1F6B1E}
\definecolor{lightgreen}{HTML}{E8F7E8}
\definecolor{darknavy}{HTML}{1A1A2E}
\definecolor{midgray}{HTML}{6B7280}
\definecolor{lightgray}{HTML}{F3F4F6}
\definecolor{bodytext}{HTML}{1F2937}
\definecolor{codegreen}{HTML}{9AE6B4}
\definecolor{warnred}{HTML}{FEE2E2}
\definecolor{warnborder}{HTML}{EF4444}

% ── Beamer theme ────────────────────────────────────────────────────────────
\usetheme{default}
\usecolortheme{default}
\usefonttheme{professionalfonts}
\setbeamertemplate{navigation symbols}{}
\setbeamertemplate{itemize items}{\textcolor{wurgreen}{\textbullet}}
\setbeamertemplate{itemize subitem}{\textcolor{midgray}{--}}
\setbeamertemplate{enumerate items}{\textcolor{wurgreen}{\textbf{\insertenumlabel.}}}
\setbeamercolor{frametitle}{fg=white, bg=darknavy}
\setbeamertemplate{frametitle}{%
  \begin{beamercolorbox}[wd=\paperwidth, ht=1.1cm, dp=0.2cm, leftskip=0.6cm]{frametitle}
    \usebeamerfont{frametitle}\insertframetitle
  \end{beamercolorbox}%
}
\setbeamerfont{frametitle}{size=\large, series=\bfseries}
\setbeamercolor{background canvas}{bg=white}
\setbeamercolor{normal text}{fg=bodytext}
\setbeamerfont{normal text}{size=\small}
\setbeamertemplate{footline}{%
  \begin{beamercolorbox}[wd=\paperwidth, ht=0.4cm, dp=0.15cm]{footline}%
    \hspace{0.5cm}%
    \textcolor{midgray}{\tiny Pixels, Frames \& Transformers · \coursemodule\ · Wageningen University \& Research}%
    \hfill%
    \textcolor{midgray}{\tiny \insertframenumber}%
    \hspace{0.4cm}%
  \end{beamercolorbox}%
}

% ── Boxes ───────────────────────────────────────────────────────────────────
\tcbset{wurbase/.style={arc=3pt, boxrule=0.6pt, left=7pt, right=7pt, top=5pt, bottom=5pt,
                        fontupper=\small, before skip=4pt, after skip=4pt}}
\newtcolorbox{greenbox}[1][]{wurbase, colback=lightgreen, colframe=wurgreen, #1}
\newtcolorbox{graybox}[1][]{wurbase, colback=lightgray, colframe=lightgray, #1}
\newtcolorbox{navybox}[1][]{wurbase, colback=darknavy, colframe=darknavy, colupper=white, #1}
\newtcolorbox{warnbox}[1][]{wurbase, colback=warnred, colframe=warnborder, #1}
\newtcolorbox{codebox}[1][]{wurbase, colback=darknavy, colframe=darknavy, colupper=white,
                            fontupper=\ttfamily\scriptsize, left=6pt, top=4pt, bottom=4pt, #1}

% ── Title slide ─────────────────────────────────────────────────────────────
\newcommand{\moduletitle}[1]{%
{
\setbeamertemplate{footline}{}
\begin{frame}[plain]
  \begin{tikzpicture}[remember picture, overlay]
    \fill[darknavy] (current page.south west) rectangle (current page.north east);
    \fill[wurgreen] (current page.south west) rectangle ([xshift=0.45cm]current page.north west);
    \node[anchor=west, text=white, font=\huge\bfseries, text width=13cm]
      at ([xshift=1.1cm, yshift=1.2cm]current page.west) {Pixels, Frames \& Transformers};
    \node[anchor=west, text=wurgreen, font=\large, text width=13cm]
      at ([xshift=1.1cm, yshift=0.05cm]current page.west) {#1 · \coursemodule};
    \node[anchor=west, text=midgray, font=\small, text width=12cm]
      at ([xshift=1.1cm, yshift=-1.5cm]current page.west)
      {Denise J. Roth, PhD Candidate\\
       Strategic Communication Group · Wageningen University \& Research};
  \end{tikzpicture}
\end{frame}
}}

% ── Small helpers ───────────────────────────────────────────────────────────
\newcommand{\keyterm}[1]{\textbf{\textcolor{darkgreen}{#1}}}
\newcommand{\muted}[1]{\textcolor{midgray}{#1}}
\newcommand{\credit}[1]{{\tiny\textcolor{midgray}{#1}}}
\newcommand{\notebookslide}[1]{%
\begin{frame}{In the Notebook}
\vspace{0.1cm}
#1
\end{frame}}
\newcommand{\nbitem}[2]{\item[\textcolor{wurgreen}{#1}] #2}

%   \begin{document}
%   \moduletitle{Images as Data}
%   ...
% Compile with: latexmk -pdf <deck>.tex
% ════════════════════════════════════════════════════════════════════════════

\documentclass[aspectratio=169]{beamer}

% ── Packages ────────────────────────────────────────────────────────────────
\usepackage[T1]{fontenc}
\usepackage[utf8]{inputenc}
\usepackage{lmodern}
\usepackage{microtype}
\usepackage{xcolor}
\usepackage{tikz}
\usetikzlibrary{arrows.meta, positioning, calc, shapes.geometric, decorations.pathreplacing, matrix, fit}
\usepackage{booktabs}
\usepackage{fontawesome5}
\usepackage{array}
\usepackage[most]{tcolorbox}

\providecommand{\coursemodule}{}

% ── WUR color palette ───────────────────────────────────────────────────────
\definecolor{wurgreen}{HTML}{34B233}
\definecolor{darkgreen}{HTML}{1F6B1E}
\definecolor{lightgreen}{HTML}{E8F7E8}
\definecolor{darknavy}{HTML}{1A1A2E}
\definecolor{midgray}{HTML}{6B7280}
\definecolor{lightgray}{HTML}{F3F4F6}
\definecolor{bodytext}{HTML}{1F2937}
\definecolor{codegreen}{HTML}{9AE6B4}
\definecolor{warnred}{HTML}{FEE2E2}
\definecolor{warnborder}{HTML}{EF4444}

% ── Beamer theme ────────────────────────────────────────────────────────────
\usetheme{default}
\usecolortheme{default}
\usefonttheme{professionalfonts}
\setbeamertemplate{navigation symbols}{}
\setbeamertemplate{itemize items}{\textcolor{wurgreen}{\textbullet}}
\setbeamertemplate{itemize subitem}{\textcolor{midgray}{--}}
\setbeamertemplate{enumerate items}{\textcolor{wurgreen}{\textbf{\insertenumlabel.}}}
\setbeamercolor{frametitle}{fg=white, bg=darknavy}
\setbeamertemplate{frametitle}{%
  \begin{beamercolorbox}[wd=\paperwidth, ht=1.1cm, dp=0.2cm, leftskip=0.6cm]{frametitle}
    \usebeamerfont{frametitle}\insertframetitle
  \end{beamercolorbox}%
}
\setbeamerfont{frametitle}{size=\large, series=\bfseries}
\setbeamercolor{background canvas}{bg=white}
\setbeamercolor{normal text}{fg=bodytext}
\setbeamerfont{normal text}{size=\small}
\setbeamertemplate{footline}{%
  \begin{beamercolorbox}[wd=\paperwidth, ht=0.4cm, dp=0.15cm]{footline}%
    \hspace{0.5cm}%
    \textcolor{midgray}{\tiny Pixels, Frames \& Transformers · \coursemodule\ · Wageningen University \& Research}%
    \hfill%
    \textcolor{midgray}{\tiny \insertframenumber}%
    \hspace{0.4cm}%
  \end{beamercolorbox}%
}

% ── Boxes ───────────────────────────────────────────────────────────────────
\tcbset{wurbase/.style={arc=3pt, boxrule=0.6pt, left=7pt, right=7pt, top=5pt, bottom=5pt,
                        fontupper=\small, before skip=4pt, after skip=4pt}}
\newtcolorbox{greenbox}[1][]{wurbase, colback=lightgreen, colframe=wurgreen, #1}
\newtcolorbox{graybox}[1][]{wurbase, colback=lightgray, colframe=lightgray, #1}
\newtcolorbox{navybox}[1][]{wurbase, colback=darknavy, colframe=darknavy, colupper=white, #1}
\newtcolorbox{warnbox}[1][]{wurbase, colback=warnred, colframe=warnborder, #1}
\newtcolorbox{codebox}[1][]{wurbase, colback=darknavy, colframe=darknavy, colupper=white,
                            fontupper=\ttfamily\scriptsize, left=6pt, top=4pt, bottom=4pt, #1}

% ── Title slide ─────────────────────────────────────────────────────────────
\newcommand{\moduletitle}[1]{%
{
\setbeamertemplate{footline}{}
\begin{frame}[plain]
  \begin{tikzpicture}[remember picture, overlay]
    \fill[darknavy] (current page.south west) rectangle (current page.north east);
    \fill[wurgreen] (current page.south west) rectangle ([xshift=0.45cm]current page.north west);
    \node[anchor=west, text=white, font=\huge\bfseries, text width=13cm]
      at ([xshift=1.1cm, yshift=1.2cm]current page.west) {Pixels, Frames \& Transformers};
    \node[anchor=west, text=wurgreen, font=\large, text width=13cm]
      at ([xshift=1.1cm, yshift=0.05cm]current page.west) {#1 · \coursemodule};
    \node[anchor=west, text=midgray, font=\small, text width=12cm]
      at ([xshift=1.1cm, yshift=-1.5cm]current page.west)
      {Denise J. Roth, PhD Candidate\\
       Strategic Communication Group · Wageningen University \& Research};
  \end{tikzpicture}
\end{frame}
}}

% ── Small helpers ───────────────────────────────────────────────────────────
\newcommand{\keyterm}[1]{\textbf{\textcolor{darkgreen}{#1}}}
\newcommand{\muted}[1]{\textcolor{midgray}{#1}}
\newcommand{\credit}[1]{{\tiny\textcolor{midgray}{#1}}}
\newcommand{\notebookslide}[1]{%
\begin{frame}{In the Notebook}
\vspace{0.1cm}
#1
\end{frame}}
\newcommand{\nbitem}[2]{\item[\textcolor{wurgreen}{#1}] #2}


\begin{document}

\moduletitle{Bias, Measurement and Faces}

% ─────────────────────────────────────────────────────────────────────────────
\begin{frame}{Errors That Aren't Random}
\begin{columns}[T]
\begin{column}{0.52\textwidth}
\textbf{Gender Shades} \muted{(Buolamwini \& Gebru, 2018)} audited commercial gender
classifiers on a set of faces balanced by skin type and gender.

\begin{itemize}
\item Two benchmarks widely used for face analysis were about 80\% and 86\% lighter-skinned
\item The classifiers' error rates ranged from under 1\% for lighter-skinned men to as much as
      35\% for darker-skinned women
\end{itemize}
\end{column}
\begin{column}{0.45\textwidth}
\begin{warnbox}
\textbf{The lesson for any vision model:} a classifier that is accurate on average can be much
worse for some groups, usually the ones that were rare in its training data.
\end{warnbox}
\begin{graybox}
\textbf{Why social scientists should care:} if errors differ between the groups you compare, your
comparison is biased (slide 8).
\end{graybox}
\end{column}
\end{columns}
\end{frame}

% ─────────────────────────────────────────────────────────────────────────────
\begin{frame}{Who Is in the Data, and Who Labeled It?}
\begin{columns}[T]
\begin{column}{0.48\textwidth}
\textbf{FairFace} \muted{(K\"arkk\"ainen \& Joo, 2021)}
\begin{itemize}
\item About 108,000 faces, mostly from public Flickr photos
\item Collected until seven race groups were roughly balanced, also by gender and age
\item Labels from \keyterm{crowd workers} who looked at each face
\end{itemize}
\muted{CC BY 4.0. Built for exactly what we need: checking performance \emph{per group}.}
\end{column}
\begin{column}{0.48\textwidth}
\begin{warnbox}
\textbf{What the labels are, and are not}\\
Gender labels are binary and record what annotators \emph{perceived}. Race labels are annotators'
guesses at one particular scheme of social categories. Neither says how a person identifies.
\end{warnbox}
\muted{In this module we train a gender classifier only to audit it. Don't use such a model to
label real people.}
\end{column}
\end{columns}
\end{frame}

% ─────────────────────────────────────────────────────────────────────────────
\begin{frame}{Audit Design: Train Skewed, Test Balanced}
\vspace{0.1cm}
\centering
\begin{tikzpicture}[
  d/.style={draw=midgray, fill=lightgray, rounded corners=3pt, font=\footnotesize, align=center,
            inner sep=5pt, text width=3.0cm},
  g/.style={draw=wurgreen, fill=lightgreen, rounded corners=3pt, font=\footnotesize, align=left,
            inner sep=5pt, text width=4.6cm},
  n/.style={draw=darknavy, fill=darknavy, text=white, rounded corners=3pt, font=\footnotesize,
            align=left, inner sep=5pt, text width=4.6cm},
  arr/.style={-{Stealth[length=2mm]}, thick, color=midgray}]
\node[d] (all) at (0,0) {FairFace validation faces, with DINOv2 embeddings};
\node[g] (pool) at (5.3,1.05) {\textbf{pool}: draw a skewed training sample (60\% White, 15\% East
Asian, \ldots, 2\% Black), half Male and half Female in every group};
\node[n] (audit) at (5.3,-1.05) {\textbf{audit half}, kept balanced, used only to measure errors};
\node[d, text width=2.5cm] (res) at (10.6,0) {error rate \textbf{per group}, with 95\% intervals};
\draw[arr] (all.east) -- (pool.west);
\draw[arr] (all.east) -- (audit.west);
\draw[arr] (pool.east) -- node[above right, font=\scriptsize, color=midgray] {train} (res.north west);
\draw[arr] (audit.east) -- node[below right, font=\scriptsize, color=midgray] {test} (res.south west);
\end{tikzpicture}

\vspace{0.35cm}
\raggedright
The skew is in \emph{who appears}, not in the outcome. Then repeat with a balanced training
sample of the same size: does the gap close? Often it narrows; it may not vanish, because the
pre-trained embeddings carry their own training data's skew.
\end{frame}

% ─────────────────────────────────────────────────────────────────────────────
\begin{frame}{Report Error Rates per Group, with Intervals}
\begin{columns}[T]
\begin{column}{0.5\textwidth}
\begin{itemize}
\item One overall accuracy hides how errors are distributed
\item Report the error rate \keyterm{per group}, and per \keyterm{intersection} (race $\times$
      gender): Gender Shades found the largest errors there
\item Use the \keyterm{bootstrap}: resample each group's faces many times and see how much its
      error rate moves
\end{itemize}
\end{column}
\begin{column}{0.47\textwidth}
\textbf{How precise is an error rate of 10\%?}

\vspace{0.1cm}
{\footnotesize
\begin{tabular}{@{}rc@{}}
\toprule
faces in the group & 95\% interval \\
\midrule
50 & $\pm$ 8.3 points \\
200 & $\pm$ 4.2 points \\
800 & $\pm$ 2.1 points \\
\bottomrule
\end{tabular}}

\vspace{0.15cm}
\muted{To tell a gap of a few points from noise, you need hundreds of faces \emph{per subgroup}.}
\end{column}
\end{columns}
\end{frame}

% ─────────────────────────────────────────────────────────────────────────────
\begin{frame}{What Do Embeddings Encode About People?}
\begin{columns}[T]
\begin{column}{0.5\textwidth}
A \keyterm{probe}: try to predict an attribute from the embeddings with cross-validation. If it
works, the information is there, and \textbf{any} model trained on the embeddings can use it.

\vspace{0.2cm}
\textbf{On the Congress portraits}, we probe for attributes from the public metadata (never
inferred from faces): party, gender, chamber, and whether a member entered Congress in 2017 or
later.
\end{column}
\begin{column}{0.47\textwidth}
\begin{graybox}
\textbf{Two kinds of information}\\
about the \emph{people} (appearance correlated with gender or age), and about the \emph{photos}
(backdrop, lighting, the year it was taken, the chamber's studio).
\end{graybox}
\begin{warnbox}
\textbf{Shortcut learning:} a model trained to code ``formal versus casual'' photos may quietly
learn gender or photo era instead. Audit it across those attributes.
\end{warnbox}
\end{column}
\end{columns}
\end{frame}

% ─────────────────────────────────────────────────────────────────────────────
\begin{frame}{Categories, Skin Tone and Identity}
\begin{columns}[T]
\begin{column}{0.31\textwidth}
\begin{graybox}[height=3.7cm]
\raggedright\textbf{Annotated categories}\\
someone's perception of a face, in someone's scheme of categories (FairFace's seven groups)
\end{graybox}
\end{column}
\begin{column}{0.31\textwidth}
\begin{greenbox}[height=3.7cm]
\raggedright\textbf{Skin tone scales}\\
appearance, not identity: Fitzpatrick (6 types) or the 10-point \textbf{Monk Skin Tone} scale
(Google, 2022), with more detail for darker tones
\end{greenbox}
\end{column}
\begin{column}{0.31\textwidth}
\begin{navybox}[height=3.7cm]
\raggedright\textbf{Self-report}\\
measures identity directly, with consent, but rarely exists for images from the web
\end{navybox}
\end{column}
\end{columns}

\vspace{0.2cm}
\begin{warnbox}
\textbf{Inferring gender, race, age or emotion from a face} is a measurement problem as much as an
ethical one. Alternatives: use metadata or self-reports; code the \emph{context} (role, framing)
instead of identity; or have trained coders code perceived attributes and report agreement.
\end{warnbox}
\end{frame}

% ─────────────────────────────────────────────────────────────────────────────
\begin{frame}{When Errors Differ by Group, Conclusions Shift}
\begin{columns}[T]
\begin{column}{0.46\textwidth}
\centering
\begin{tikzpicture}[x=1cm, y=4.5cm]
\draw[midgray] (0,0) -- (4.6,0);
\foreach \x/\h/\c/\l in {0.5/0.40/midgray/40, 1.3/0.42/wurgreen/42, 2.7/0.40/midgray/40, 3.5/0.34/wurgreen/34}
  { \fill[\c] (\x-0.3,0) rectangle (\x+0.3,\h); \node[above, font=\scriptsize] at (\x,\h) {\l\%}; }
\node[below, font=\footnotesize] at (0.9,0) {group A};
\node[below, font=\footnotesize] at (3.1,0) {group B};
\fill[midgray] (0.2,0.62) rectangle (0.45,0.66); \node[right, font=\scriptsize] at (0.5,0.64) {true share smiling};
\fill[wurgreen] (0.2,0.55) rectangle (0.45,0.59); \node[right, font=\scriptsize] at (0.5,0.57) {measured by the detector};
\end{tikzpicture}\\
\muted{\footnotesize expected values in the simulation}
\end{column}
\begin{column}{0.51\textwidth}
A smile detector finds 90\% of smiles in group A but 70\% in group B, with 10\% false alarms in
both. It is right 86\% of the time, and it \textbf{invents} an 8-point gap.

\vspace{0.15cm}
\begin{greenbox}
\textbf{Fix: a validation sample per group.} Hand-code a random sample, estimate each group's
sensitivity and false positive rate, and correct \muted{(Rogan \& Gladen, 1978)}:\\[3pt]
{\footnotesize $\text{corrected} = \dfrac{\text{measured} - \text{false positive rate}}{\text{sensitivity} - \text{false positive rate}}$}
\end{greenbox}
\muted{\footnotesize Regression methods that carry this uncertainty through:
TeBlunthuis et al.\ (2024).}
\end{column}
\end{columns}
\end{frame}

% ─────────────────────────────────────────────────────────────────────────────
\begin{frame}{What to Report}
When a measurement in your study involves people in images, readers need to know:

\vspace{0.2cm}
\begin{enumerate}
\setlength{\itemsep}{4pt}
\item \textbf{Who is in your data}, and how that compares with the population you make claims
      about
\item \textbf{How labels about people were obtained}: metadata, self-report, or someone's
      perception, and in which categories
\item \textbf{Error rates per relevant group}, with intervals and the size of each group in your
      validation sample
\item \textbf{What you did about differences}: rebalancing, correction, or a stated limitation
\item \textbf{What you chose not to infer}, and why
\end{enumerate}
\end{frame}

% ─────────────────────────────────────────────────────────────────────────────
\notebookslide{
\begin{itemize}
\setlength{\itemsep}{5pt}
\nbitem{\faUsers}{FairFace: who is in it, and how it was labeled}
\nbitem{\faBalanceScale}{A Gender Shades-style audit: skewed versus balanced training data}
\nbitem{\faChartBar}{Error rates per group and per intersection, with bootstrap intervals}
\nbitem{\faSearch}{Probing portrait embeddings for people and for photo production}
\nbitem{\faPalette}{Skin tone, categories and identity}
\nbitem{\faCalculator}{A simulation of differential error, and a correction with a validation sample}
\end{itemize}

\vspace{0.35cm}
\muted{\footnotesize part2\_vectors/07\_bias\_measurement\_and\_faces/notebooks/\\
07\_bias\_measurement\_and\_faces\_exercises.ipynb · 5 exercises · GPU recommended}
}

\end{document}
